\documentclass{article}
\usepackage[T1]{fontenc}
\usepackage{lmodern}
\usepackage{graphicx}
\usepackage{amsmath,amssymb}
\usepackage[a4paper,margin=2cm]{geometry}
\usepackage[absolute,overlay]{textpos}
\setlength{\TPHorizModule}{1mm}
\setlength{\TPVertModule}{1mm}
\usepackage[indonesian]{babel}
\usepackage{hyperref}
\setlength{\emergencystretch}{2em}
% unit: Halaman 33

\begin{document}

% === Halaman 33 (python/native) ===

Twofish

• Salah satu finalis AES (Advanced Encryption Standard)

• Dirancang oleh Bruce Schneier.

• Twofish dapat dianggap sebagai kelanjutan dari Blowfish

• Ukuran blok plainteks dan cipherteks di dalam Twofish adalah 128 bit

• Panjang kunci sampai 256 bit

• Jumlah putaran 16 kali, setiap putaran melakukan operasi substitusi-

permutasi

• Jumlah kotak S-box adalah 5 buah 

33

\end{document}
